% =====================================================================
%  Seminar 1 — Plant-fungal relationships and their impact on the ecosystem
% =====================================================================
\section{Plant-fungal relationships and their impact on the ecosystem}

\subsection{About the webinar}

\begin{infotable}
Title & Plant-fungal relationships and their impact on the ecosystem \row
Link & \href{https://bsbi.org/take-part/events/calendar/webinar-plant-fungal-relationships-and-their-impact-on-the-ecosystem}{bsbi.org — event page} \row
Summary & ``As we move into autumn, the vital connections between plants and fungi become more visible than ever. Fungi play many important roles in the environment, including forming relationships with plant partners. Mycorrhizal fungi exchange nutrients and water with different grasses, trees, and even flowers. Learn more about how these relationships work and what mycorrhizal fungi to spot in Ireland. There will be a short introduction to the biology of these fungi, followed by their habitats and plant partners. Finally, we will discuss their biggest threats and how you can support the plant-fungal relationships in your garden.'' \row
Speaker & Delia Termini, MSc Environmental Science (Trinity College Dublin) \row
Host & Kim Lake, for BSBI's Botanical Skills Northern Ireland project \row
Date and time & 13.10.2026; 19:00--20:00 BST (20:00--21:00 CET) \row
Location & Online webinar (Zoom), with audience Q\&A; recorded and uploaded to the BSBI YouTube channel \row
Portfolio & Seminar 1 on the Seminar Signature Sheet; digitally signed by the speaker
\end{infotable}

\subsubsection{Website}
\begin{center}\img{images/seminar01/website.jpg}\end{center}

% ---------------------------------------------------------------------
\subsection{Registration and proof of attendance}
\begin{center}
\img{images/seminar01/registration.jpg}\\[2pt]
{\footnotesize\color{muted}Registration}
\end{center}

% The speaker's e-signed attendance confirmation (PDF):
\attachpdf{Attendance confirmation e-signed by the speaker}{pdfs/seminar01-attendance.pdf}{subsection}

% ---------------------------------------------------------------------
\subsection{Notes}
Session funded by the Department of Agriculture, Environment and Rural Affairs (DAERA) via the Carrier Bag Levy. Delia Termini completed an MSc in Environmental Science focused on fungal specimens in fungaria in Ireland.

\begin{itemize}
  \item Research: comparing fungarium records (from as early as 1806; 5--209 per county) with citizen science (iNaturalist; 41--2,065 per county) to identify correlations in fungal distributions across Ireland
  \item Recording effort follows where collectors and app users are, not only where fungi are $\rightarrow$ distribution data carry observer bias
\end{itemize}

\begin{imagerow}
  \rowimg{0.48}{images/seminar01/fungaria-map.jpg}{Fungaria specimen richness} \hfill
  \rowimg{0.48}{images/seminar01/inaturalist-map.jpg}{iNaturalist specimen richness}
\end{imagerow}

% .....................................................................
\subsubsection{Fungal biodiversity}

\begin{imagerow}
  \rowimg{0.22}{images/seminar01/tar-spot.jpg}{} \hfill
  \rowimg{0.22}{images/seminar01/bracket.jpg}{} \hfill
  \begin{minipage}[t]{0.22\linewidth}\centering
    \img{images/seminar01/crust-1.jpg}\\[6pt]\img{images/seminar01/crust-2.jpg}
  \end{minipage} \hfill
  \rowimg{0.22}{images/seminar01/coral-fungus.jpg}{}
\end{imagerow}

\begin{itemize}
  \item Fungi are their own kingdom
  \item The mushroom is the fruiting body; the organism is the mycelium (network of hyphae)
\end{itemize}

\begin{imagerow}
  \rowimg{0.48}{images/seminar01/mycelium-log.jpg}{} \hfill
  \rowimg{0.48}{images/seminar01/mushrooms-cone.jpg}{}
\end{imagerow}

\begin{itemize}
  \item Biodiversity covers variation at genetic, species, and ecosystem levels
  \item Fungi are at the base of the soil food web as decomposers and symbiotic partners to plants
\end{itemize}

\begin{imagerow}
  \rowimg{0.48}{images/seminar01/soil-food-web.jpg}{} \hfill
  \rowimg{0.48}{images/seminar01/plant-network.jpg}{}
\end{imagerow}

\begin{itemize}
  \item Estimated vs.\ described fungal species (adapted from Niskanen \textit{et al.}, 2023): 2.5 million vs.\ 155,000
  \begin{itemize}
    \item Compared to plants (89\% described), fungi lag far behind (\textasciitilde6\%)
  \end{itemize}
  \item Largest organism in the world: a dark honey fungus (\textit{Armillaria ostoyae}) in Oregon = one genetic individual spread across a whole forest (\textasciitilde5\,mi\textsuperscript{2})
\end{itemize}

\begin{imagerow}
  \rowimg{0.48}{images/seminar01/armillaria-map.jpg}{} \hfill
  \rowimg{0.48}{images/seminar01/honey-fungus.jpg}{}
\end{imagerow}

\sidebyside{0.45}{%
  \begin{itemize}
    \item Global fungal diversity hotspots
  \end{itemize}}{images/seminar01/hotspots.jpg}

% .....................................................................
\subsubsection{Fungal biology}
\paragraph{Identification}\mbox{}

\begin{center}\img[0.55\linewidth]{images/seminar01/amanita-diagram.jpg}\end{center}

\begin{enumerate}
  \item Cap colour, shape, texture (over time):
\end{enumerate}
\begin{imagerow}
  \rowimg{0.3}{images/seminar01/cap-1.jpg}{} \hfill
  \rowimg{0.3}{images/seminar01/cap-2.jpg}{} \hfill
  \rowimg{0.3}{images/seminar01/cap-3.jpg}{}
\end{imagerow}

\begin{enumerate}[resume]
  \item Spore-bearing surface (hymenium):
\end{enumerate}
\begin{imagerow}
  \rowimg{0.2}{images/seminar01/spines.jpg}{Spines (very rare)} \hfill
  \rowimg{0.36}{images/seminar01/gills.jpg}{Gills} \hfill
  \rowimg{0.36}{images/seminar01/pores.jpg}{Pores}
\end{imagerow}

\begin{enumerate}[resume]
  \item Spores (colour, microscopy, spore prints):
  \begin{enumerate}[label=\alph*.]
    \item Dispersal methods vary vastly between species: wind from gills; puffballs released by raindrops/touch; cup fungi shoot spores out; bird's nest fungi splashed by rain; stinkhorns spread by insects (attracted to the smell); inkcaps dissolve (autodigest) to release spores
  \end{enumerate}
\end{enumerate}
\begin{imagerow}
  \rowimg{0.3}{images/seminar01/inkcap.jpg}{} \hfill
  \rowimg{0.3}{images/seminar01/shaggy-inkcap.jpg}{} \hfill
  \rowimg{0.3}{images/seminar01/birds-nest.jpg}{}
\end{imagerow}

\begin{enumerate}[resume]
  \item Stem (stipe): shape, colour
  \item Ring (annulus): presence, shape, texture
  \item Flesh: texture, colour, colour change, smell, milk (latex)
  \item Stem base: shape, volva, mycelia
  \item Substrate: what is it growing on/with/in?
\end{enumerate}

\sidebyside{0.45}{%
  \textbf{Fungal lifecycle}\par
  Spores germinate into hyphae $\rightarrow$ hyphae of two compatible mating types fuse $\rightarrow$ mated mycelium grows and forms primordia $\rightarrow$ primordia develop into fruiting bodies $\rightarrow$ these release new spores.}%
  {images/seminar01/lifecycle.jpg}

% .....................................................................
\subsubsection{Saprotrophic fungi}

\begin{imagerow}
  \rowimg{0.26}{images/seminar01/sapro-1.jpg}{} \hfill
  \begin{minipage}[t]{0.2\linewidth}\centering
    \img{images/seminar01/sapro-2.jpg}\\[6pt]\img{images/seminar01/sapro-3.jpg}
  \end{minipage} \hfill
  \rowimg{0.46}{images/seminar01/mycelium-soil.jpg}{}
\end{imagerow}

\sidebyside{0.45}{%
  Decomposers: draw nutrients from decaying matter via their mycelium through extracellular digestion (enzyme secretion) $\rightarrow$ very important to the plants that have a symbiotic relationship with them}%
  {images/seminar01/enzymes.jpg}

\begin{itemize}
  \item Break down dead wood and leaf litter, making C and nutrients available for uptake again
  \item Without these fungi, dead plant matter would build up
  \item Global importance of fungi
  \begin{itemize}
    \item The breakdown releases CO\textsubscript{2} (more fungal competition during breakdown = slower CO\textsubscript{2} release $\rightarrow$ studies show that communities with higher fungal diversity release less CO\textsubscript{2})
  \end{itemize}
  \item Can also be parasitic:
  \begin{itemize}
    \item Plant-specific
    \item Can limit seeds
    \item Parasitic fungi mostly attack weak members/organs $\rightarrow$ allows other plants to compete
    \begin{itemize}
      \item E.g.\ honey fungus: helps kill/break down older and weaker trees, letting younger trees take their place in the new canopy opening
    \end{itemize}
  \end{itemize}
\end{itemize}

\needspace{12\baselineskip}
\paragraph{Plant associations with saprotrophic fungi in Ireland}\mbox{}

% A table replaces Notion's five squashed columns
\begin{center}
\arrayrulecolor{rule}\renewcommand{\arraystretch}{1.4}
\begin{tabularx}{\linewidth}{|>{\centering\arraybackslash}m{2.6cm}|>{\raggedright\arraybackslash}p{3.6cm}|>{\raggedright\arraybackslash}X|}
\hline
\rowcolor{shade} & \textbf{Species} & \textbf{Notes} \\ \hline
\img[2.4cm]{images/seminar01/scarlet-elf-cup-1.jpg}\par\vspace{3pt}\img[2.4cm]{images/seminar01/scarlet-elf-cup-2.jpg}
  & Scarlet elf cup \newline (\textit{Sarcoscypha austriaca} / \textit{S.\ coccinea}; indistinguishable without microscopy)
  & \tocheck{Grows on fallen, mossy twigs; does not grow with pine; needs very specific conditions to grow (moist, with damp wood)} \\ \hline
\img[2.4cm]{images/seminar01/earthstar.jpg}
  & Collared earthstar \newline (\textit{Geastrum triplex})
  & Only able to colonise already decaying matter (does not grow on logs) \\ \hline
\img[2.4cm]{images/seminar01/dryads-saddle.jpg}
  & Dryad's saddle \newline (\textit{Polyporus squamosus})
  & Somewhat parasitic: will take over a not fully dead tree (e.g.\ a missing limb) $\rightarrow$ negative impact; smells like watermelon \\ \hline
\img[2.4cm]{images/seminar01/porcelain.jpg}
  & Porcelain fungus \newline (\textit{Mucidula mucida})
  & Only grows on beech trees \\ \hline
\img[2.4cm]{images/seminar01/birch-polypore.jpg}
  & Birch polypore \newline (\textit{Fomitopsis betulina})
  & Only grows on birch trees \\ \hline
\end{tabularx}
\end{center}

% .....................................................................
\subsubsection{Mycorrhizal fungi}

\sidebyside{0.45}{%
  \begin{itemize}
    \item Colonise living roots $\rightarrow$ exchange of nutrients and water from the fungi for sugars from the plant
    \begin{itemize}
      \item Also add to soil organic C through hyphal biomass and soil aggregation
      \item Some can prevent plant heavy-metal uptake $\rightarrow$ improve chemical and biological soil properties and increase plant survival in stressful environments
    \end{itemize}
  \end{itemize}}%
  {images/seminar01/mycorrhiza-overview.jpg}

\begin{itemize}
  \item Arbuscular mycorrhiza (332 known species)
\end{itemize}
\begin{imagerow}
  \rowimg{0.48}{images/seminar01/am-diagram.jpg}{} \hfill
  \rowimg{0.48}{images/seminar01/am-microscopy.jpg}{}
\end{imagerow}

\begin{itemize}
  \item Ectomycorrhizal fungi (do not enter the cell; \textasciitilde5,000 species known/estimated) $\rightarrow$ mostly without fruiting bodies
  \begin{itemize}
    \item Orchids $\rightarrow$ more mycelium = more C uptake
  \end{itemize}
\end{itemize}
\begin{imagerow}
  \rowimg{0.24}{images/seminar01/ecm-diagram.jpg}{} \hfill
  \rowimg{0.38}{images/seminar01/ecm-section.jpg}{} \hfill
  \rowimg{0.3}{images/seminar01/orchid-carbon.jpg}{}
\end{imagerow}

\begin{itemize}
  \item Ericoid mycorrhizal fungi (only associated with Ericaceae)
\end{itemize}
\begin{imagerow}
  \rowimg{0.45}{images/seminar01/ericoid-microscopy.jpg}{} \hfill
  \rowimg{0.4}{images/seminar01/ericoid-drawing.jpg}{}
\end{imagerow}

\textbf{Amanitaceae:} associated with birch, beech, pines, and oaks, and occasionally spruce and willow.
\begin{imagerow}
  \rowimg{0.3}{images/seminar01/amanita-1.jpg}{} \hfill
  \rowimg{0.3}{images/seminar01/amanita-2.jpg}{} \hfill
  \rowimg{0.3}{images/seminar01/amanita-3.jpg}{}
\end{imagerow}

\textbf{Russulaceae:} typically prefer conifers but can grow with beeches.
\begin{itemize}
  \item Milk caps: exude latex and are frequently found with birch; can associate with pines and beeches
\end{itemize}
\begin{imagerow}
  \rowimg{0.3}{images/seminar01/russula-1.jpg}{} \hfill
  \rowimg{0.3}{images/seminar01/russula-2.jpg}{} \hfill
  \rowimg{0.3}{images/seminar01/russula-3.jpg}{}
\end{imagerow}

\textbf{Boletaceae:} have pores and a fatter, squishy stipe; specific species are associated with birch, oak, hazel, and spruce.
\begin{imagerow}
  \rowimg{0.3}{images/seminar01/bolete-1.jpg}{} \hfill
  \rowimg{0.3}{images/seminar01/bolete-2.jpg}{} \hfill
  \rowimg{0.3}{images/seminar01/bolete-3.jpg}{}
\end{imagerow}

% ---------------------------------------------------------------------
\subsection{Threats}

\begin{itemize}
  \item Climate change (droughts)
  \item Invasive species (e.g.\ the Asian \textit{Favolaschia calocera}) $\rightarrow$ threaten Irish and UK native saprotrophic species because they are better suited to the warming conditions and have no natural predators
\end{itemize}
\begin{imagerow}
  \rowimg{0.3}{images/seminar01/favolaschia.jpg}{} \hfill
  \rowimg{0.45}{images/seminar01/oyster-grey.jpg}{}
\end{imagerow}

\sidebyside{0.45}{%
  \begin{itemize}
    \item Agriculture and development (monoculture; compaction; chemicals and fungicides $\rightarrow$ impact the soil biome)
  \end{itemize}}%
  {images/seminar01/soil-cores.jpg}

\textbf{Waxcaps = threatened}
\begin{itemize}
  \item Only grow on unimproved grasslands
  \item Ballerina waxcap is on the IUCN Red List
  \begin{itemize}
    \item Current study in Ireland: recording the number of species for conservation purposes
  \end{itemize}
\end{itemize}

\subsubsection{How can we support fungi in our environment?}
\begin{itemize}
  \item Grassland fungi: unimproved soil; remove dead matter; shorter grass (not dead) $\rightarrow$ do best with high plant and fungal diversity (grasses, shrubs, plantains, \ldots)
  \item Mycorrhizal fungi: need certain trees; increase tree diversity, no tillage, lower fertiliser and pesticide use
  \item Saprotrophic fungi: need plant matter to decay (leave dead plant matter, e.g.\ leaves, over the winter; pile up dead logs/twigs)
\end{itemize}
Overall: leave the garden a bit `dirty' and `messy' to promote fungal growth and diversity.

% ---------------------------------------------------------------------
\subsection{Reflection}
\begin{callout}
This talk deepened my interest in fungal biology, which began last academic year with the lecture on mycorrhiza in the \textit{Terrestrial Ecology} course. The example of waxcaps in Ireland threatened by grassland management is particularly relevant to my thesis on dike grassland management, as we have not yet considered the implications for fungi.

Additionally, the use of citizen science in Delia's MSc thesis and the aim of the talk at a general audience are both particularly relevant for this course and for the course \textit{Science Communication and Outreach}, which I will take next semester.

Questions to look into: How can I integrate mycorrhizal community analysis into the thesis work with the data we have collected so far? What are the implications of the way INBO is managing and studying the dike sites in the Scheldt basin for fungal diversity and survival?
\end{callout}
